% NOETHER PAPER 6 — KOREAN TRANSLATION-PRODUCER DRAFT — UNCHECKED
% Unit: T01-U05; current-authority whole lines 4610--4612
% Source slice: 1,453 LF UTF-8 bytes; SHA-256 DDF5A0803E76C56A322378C0962A1D371836B6A1022B114015BF3B313A7D595C
% Pointer: NOETH-DE-AUTH-v006-20260804; SHA-256 DB99DD87100654674D7ED24B4ABBBBC3A9920CCF035740D276CE8A87A5313C18
% German authority: NOETH-DE-ED-0001; 2,153,565 raw UTF-8 bytes; SHA-256 D1F06B311F6CBD991DD247D745DD9A72DDE326A20396DF43CFE0C8EDB1593CDB
% This is producer translation only: no source/Korean/formula review, compilation, or rendering.

본 연구의 계기는 E. Fischer와 나눈 대화, 특히 그가 라그랑주 종영역(Gattungsbereich)의 최소기저에 관하여 제기한 질문이었다. 그와 관련된 연구는 지정된 군을 갖는 방정식의 구성으로 이어졌으며(앞에서 인용한 \glqq{}유리 함수체\grqq{}에 관한 보고의 제2부를 참조하라), 추후 발표로 남겨 둔다.

본 논문이 최초의 보고에 비해 갖는 본질적인 확장은, 그곳에서는 체나 특수한 정역에 대해서만 진술되었던 기저정리들을 임의의 계로 옮겼다는 데 있다. 이러한 이행은 임의의 계를 포함하는 최소체라는 개념, 곧 정역의 분수체를 일반화한 개념에 힘입어 간단하게 이루어진다(\S\ 7). 체의 유리기저가 알려진 뒤 K. Hentzelt가 힐베르트의 가군 기저를 거치는 우회로를 통해 임의의 계에도 유리기저가 존재함을 증명할 수 있었던 사실이 이 개념을 세우도록 나를 이끌었다. 또한 내가 정역에 적용한 추론들이 동일한 전제를 만족하는 임의의 계에도 유효하다는 K. Hentzelt의 지적은 정칙계의 정수성 기저에 관한 정리로도 나를 이끌었다. 그 밖의 몇 가지 개별적인 작은 지적 역시 Hentzelt에게 빚지고 있다.
